%\documentclass[a4paper,dvipsnames]{article}
%
%\addtolength{\hoffset}{-2.25cm}
%\addtolength{\textwidth}{4.5cm}
%\addtolength{\voffset}{-3.25cm}
%\addtolength{\textheight}{5cm}
%\setlength{\parskip}{0pt}
%\setlength{\parindent}{0in}
%
%%----------------------------------------------------------------------------------------
%	PACKAGES AND OTHER DOCUMENT CONFIGURATIONS
%----------------------------------------------------------------------------------------

%----------------------------------------------------------------------------------------
%		Generals
%----------------------------------------------------------------------------------------
\usepackage{fourier}
\usepackage{frcursive}
\usepackage[T1]{fontenc} %Accents handling
\usepackage[utf8]{inputenc} % Use UTF-8 encoding
%\usepackage{microtype} % Slightly tweak font spacing for aesthetics
\usepackage[english, francais]{babel} % Language hyphenation and typographical rules

%----------------------------------------------------------------------------------------
%		Graphics
%----------------------------------------------------------------------------------------
\usepackage{xcolor}
\usepackage{graphicx, multicol} % Enhanced support for graphics
\graphicspath{{FIG/}}
\usepackage{wrapfig}

%----------------------------------------------------------------------------------------
%		Other packages
%----------------------------------------------------------------------------------------
\usepackage{hyperref}
\hypersetup{
	colorlinks=true, %colorise les liens
	breaklinks=true, %permet le retour à la ligne dans les liens trop longs
	urlcolor= bleu3,  %couleur des hyperliens
	linkcolor= bleu3, %couleur des liens internes
	plainpages=false  %pour palier à "Bookmark problems can occur when you have duplicate page numbers, for example, if you have a page i and a page 1."
}
\usepackage{tabularx}
\newcolumntype{M}[1]{>{\arraybackslash}m{#1}} %Defines a scalable column type in tabular
\usepackage{booktabs} % Enhances quality of tables
\usepackage{diagbox} % barre en diagonale dans un tableau
\usepackage{multicol}
\usepackage[explicit]{titlesec}


%----------------------------------------------------------------------------------------
%		Headers and footers
%----------------------------------------------------------------------------------------
\usepackage{fancyhdr} % Headers and footers
\pagestyle{fancy} % All pages have headers and footers
\fancyhead{}\renewcommand{\headrulewidth}{0pt} % Blank out the default header
\renewcommand{\footrulewidth}{0pt}
\fancyfoot[L]{} % Custom footer text
\fancyfoot[C]{\href{https://sacado.xyz/}{sacado.xyz}} % Custom footer text
\fancyfoot[R]{\thepage} % Custom footer text

%----------------------------------------------------------------------------------------
%		Mathematics packages
%----------------------------------------------------------------------------------------
\usepackage{amsthm, amsmath, amssymb} % Mathematical typesetting
\usepackage{marvosym, wasysym} % More symbols
\usepackage[makeroom]{cancel}
\usepackage{xlop}
\usepackage{pgf,tikz,pgfplots}
\pgfplotsset{compat=1.15}
\usetikzlibrary{positioning}
%\usetikzlibrary{arrows}
\usepackage{pst-plot,pst-tree,pst-func, pstricks-add,pst-node,pst-text}
\usepackage{units}
\usepackage{nicefrac}
\usepackage[np]{numprint} %Séparation milliers dans un nombre

%----------------------------------------------------------------------------------------
%		New text commands
%----------------------------------------------------------------------------------------
\usepackage{calc}
\usepackage{boites}
 \renewcommand{\arraystretch}{1.6}

%%%%% Pour les imports.
\usepackage{import}

%%%%% Pour faire des boites
\usepackage[tikz]{bclogo}
\usepackage{bclogo}
\usepackage{framed}
\usepackage[skins]{tcolorbox}
\tcbuselibrary{breakable}
\tcbuselibrary{skins}
\usetikzlibrary{babel,arrows,shadows,decorations.pathmorphing,decorations.markings,patterns}

%%%%% Pour les symboles et les ensembles
\newcommand{\pp}{\leq}
\newcommand{\pg}{\geq}
%\newcommand{\euro}{\eurologo{}}
\newcommand{\R}{\mathbb{R}}
\newcommand{\N}{\mathbb{N}}
\newcommand{\D}{\mathbb{D}}
\newcommand{\Z}{\mathbb{Z}}
\newcommand{\Q}{\mathbb{Q}}
\newcommand{\C}{\mathbb{C}}

%%%%% Pour une double minipage
\newcommand{\mini}[2]{
\begin{minipage}[t]{0.48\linewidth}
#1
\end{minipage}
\hfill
\begin{minipage}[t]{0.48\linewidth}
#2
\end{minipage}
}


%\newcommand\hole[1]{\texttt{\_}}
%\newcommand{\PROP}[1]{\textbf{\underline{#1}}}
%\newcommand{\exercice}{\textcolor{OliveGreen}{Exercice : }}
%\newcommand{\correction}{\textcolor{BurntOrange}{Correction : }}
%\newcommand{\propriete}{\textbf{\underline{Propriété}} : }
%\newcommand{\prop}{\textbf{\underline{Propriété}} : }
%\newcommand{\vocabulaire}{\textbf{\underline{Vocabulaire}} : }
%\newcommand{\voca}{\textbf{\underline{Vocabulaire}} : }

\usepackage{enumitem}
\newlist{todolist}{itemize}{2} %Pour faire des QCM
\setlist[todolist]{label=$\square$} %Pour faire des QCM \begin{todolist} instead of itemize

%----------------------------------------------------------------------------------------
%		Définition de couleur pour geogebra
%----------------------------------------------------------------------------------------
\definecolor{zzttqq}{rgb}{0.6,0.2,0.} %rouge des polygones
\definecolor{qqqqff}{rgb}{0.,0.,1.}
\definecolor{xdxdff}{rgb}{0.49019607843137253,0.49019607843137253,1.}%bleu
\definecolor{qqwuqq}{rgb}{0.,0.39215686274509803,0.} %vert des angles
\definecolor{ffqqqq}{rgb}{1.,0.,0.} %rouge vif
\definecolor{uuuuuu}{rgb}{0.26666666666666666,0.26666666666666666,0.26666666666666666}
\definecolor{qqzzqq}{rgb}{0.,0.6,0.}
\definecolor{cqcqcq}{rgb}{0.7529411764705882,0.7529411764705882,0.7529411764705882} %gris
\definecolor{qqffqq}{rgb}{0.,1.,0.}
\definecolor{ffdxqq}{rgb}{1.,0.8431372549019608,0.}
\definecolor{ffffff}{rgb}{1.,1.,1.}
\definecolor{ududff}{rgb}{0.30196078431372547,0.30196078431372547,1.}

%-------------------------------------------------
%
%	EN TETE
%
%-------------------------------------------------

% Classe
\newcommand{\myClasse}   
{
    6e
}

% Discipline
\newcommand{\myDiscipline}   
{
    Mathématiques
}

% Parcours
\newcommand{\myParcours}
{
  Nombres et Calculs
}

%Titre de la séquence
\newcommand{\myTitle}
{
    \scshape\huge
\textcolor{sacado_purple}{
		Nombres décimaux
}
}

%----------------------------------------------------------------------------------------

%%----------------------------------------------------------------------------------------
%		Définition de couleur pour les boites
%----------------------------------------------------------------------------------------
\definecolor{bleu1}{rgb}{0.54,0.79,0.95} %% Bleu
\definecolor{sapgreen}{rgb}{0.4, 0.49, 0}
\definecolor{dvzfxr}{rgb}{0.7,0.4,0.}
\definecolor{beamer}{rgb}{0.5176470588235295,0.49019607843137253,0.32941176470588235} % couleur beamer
\definecolor{preuveRbeamer}{rgb}{0.8,0.4,0}
\definecolor{sectioncolor}{rgb}{0.24,0.21,0.44}
\definecolor{subsectioncolor}{rgb}{0.1,0.21,0.61}
\definecolor{subsubsectioncolor}{rgb}{0.1,0.21,0.61}
\definecolor{info}{rgb}{0.82,0.62,0}
\definecolor{bleu2}{rgb}{0.38,0.56,0.68}
\definecolor{bleu3}{rgb}{0.24,0.34,0.40}
\definecolor{bleu4}{rgb}{0.12,0.20,0.25}
\definecolor{vert}{rgb}{0.21,0.33,0}
\definecolor{vertS}{rgb}{0.05,0.6,0.42}
\definecolor{red}{rgb}{0.78,0,0}
\definecolor{color5}{rgb}{0,0.4,0.58}
\definecolor{eduscol4B}{rgb}{0.19,0.53,0.64}
\definecolor{eduscol4P}{rgb}{0.62,0.12,0.39}

%----------------------------------------------------------------------------------------
%		Définition de couleur pour les boites SACADO
%----------------------------------------------------------------------------------------
\definecolor{sacado_blue}{RGB}{0,129,159} %% Bleu Sacado
\definecolor{sacado_green}{RGB}{59, 157, 38} %% Vert Sacado
\definecolor{sacado_yellow}{RGB}{255,180,0} %% Jaune Sacado
\definecolor{sacado_purple}{RGB}{94,68,145} %% Violet foncé Sacado
\definecolor{sacado_violet}{RGB}{153,117,224} %% Violet clair Sacado
\definecolor{sacado_orange}{RGB}{249,168,100} %% Orange Sacado
\definecolor{ill_frame}{HTML}{F0F0F0}
\definecolor{ill_back}{HTML}{F7F7F7}
\definecolor{ill_title}{HTML}{AAAAAA}


 % Compteurs pour Théorème, Définition, Exemple, Remarque, .....
\newcounter{cpttheo}
\setcounter{cpttheo}{0}
\newcounter{cptdef}
\setcounter{cptdef}{0}
\newcounter{cptmth}
\setcounter{cptmth}{0}
\newcounter{cpttitre}
\setcounter{cpttitre}{0}
 % Exercices
\newcounter{cptapp}
\setcounter{cptapp}{0}
\newcounter{cptex}
\setcounter{cptex}{0}
\newcounter{cptsr}
\setcounter{cptsr}{0}
\newcounter{cpti}
\setcounter{cpti}{0}
\newcounter{cptcor}
\setcounter{cptcor}{0}




%%%%% Pour réinitialiser numéros des paragraphes après une nouvelle partie
\makeatletter
    \@addtoreset{paragraph}{part}
\makeatother


%%%% Titres et sections

\newlength\chapnumb
\setlength\chapnumb{3cm}


% \titleformat{\part}[block] {
 % \normalfont\sffamily\color{violet}}{}{0pt} {
   % \parbox[t]{\chapnumb}{\fontsize{120}{110}\selectfont\ding{110}}
   % \parbox[b]{\dimexpr\textwidth-\chapnumb\relax}{
       % \raggedleft
       % \hfill{{\color{bleu3}\fontsize{40}{30}\selectfont#1}}\\
       % \rule{0.99\textwidth-\chapnumb\relax}{0.4pt}
 % }
% }

% \titleformat{name=\part,numberless}[block]
% {\normalfont\sffamily\color{bleu3}}{}{0pt}
% {\parbox[b]{\chapnumb}{%
  % \mbox{}}%
 % \parbox[b]{\dimexpr\textwidth-\chapnumb\relax}{%
   % \raggedleft%
   % \hfill{{\color{bleu3}\fontsize{40}{30}\selectfont#1}}\\
   % \rule{0.99\textwidth-\chapnumb\relax}{0.4pt}
 % }
% }



% \titleformat{\chapter}[block] {
 % \normalfont\sffamily\color{violet}}{}{0pt} {
   % \parbox[t]{\chapnumb}{ 
     % \fontsize{120}{110}\selectfont\thechapter}
     % \parbox[b]{\textwidth-\chapnumb}{
       % \raggedleft
       % \hfill{{\color{bleu3}\huge#1}}\\  
  % \ifthenelse{\thechapter<10}{\rule{0.99\textwidth-\chapnumb}{0.4pt}}{\rule{0.9\textwidth - \chapnumb}{0.4pt}}
       % \setcounter{cpttitre}{0}
	% \setcounter{cptapp}{0}
	% \setcounter{cptex}{0}
	% \setcounter{cptsr}{0}
	% \setcounter{cpti}{0}
	% \setcounter{cptcor}{0} 
 % }
% }

% \titleformat{name=\chapter,numberless}[block]
% {\normalfont\sffamily\color{bleu3}}{}{0pt}
% {\parbox[b]{\chapnumb}{%
  % \mbox{}}%
 % \parbox[b]{\textwidth-\chapnumb}{%
   % \raggedleft
   % \hfill{{\color{bleu3}\huge#1}}\\
   % \ifthenelse{\thechapter<10}{\rule{0.99\textwidth-\chapnumb}{0.4pt}}{ \rule{0.9\textwidth - \chapnumb}{0.4pt}}
       % \setcounter{cpttitre}{0}
	% \setcounter{cptapp}{0}
	% \setcounter{cptex}{0}
	% \setcounter{cptsr}{0}
	% \setcounter{cpti}{0}
	% \setcounter{cptcor}{0} 
 % }
% }
%
%       
%
%%%%% Personnalisation des numéros des sections
\renewcommand\thesection{\Roman{section}. }
\renewcommand\thesubsection{\hspace{1cm}\arabic{subsection}. }
\renewcommand\thesubsubsection{\hspace{2cm}\alph{subsubsection}. }

\titleformat{\section}[hang]{\color{sacado_purple}{}\normalfont\filright\huge}{}{0.4em}{\textbf{\thesection  #1}}   
% \titlespacing*{\section}{0.2pt}{0ex plus 0ex minus 0ex}{0.3em}
   
\titleformat{\subsection}[hang]{\color{sacado_purple}{}\normalfont\filright\Large}{}{0.4em}{\thesubsection
 #1}            
\titleformat{\subsubsection}[hang]{\color{sacado_purple}{}\normalfont\filright\large}{}{0.4em}{\thesubsubsection
 #1}
\titleformat{\paragraph}[hang]{\color{black}{}\normalfont\filright\normalsize}{}{0.4em}{#1}



%%%%%%%%%%%%%%%%%%%%% Cycle 4
%\newcommand{\Titre}[2]{\section*{#1 
%\ifthenelse{\equal{#2}{1}}   {\hfill{ \ding{182}  \ding{173} \ding{174} } \addcontentsline{toc}{section}{#1 \ding{182}} }%
%{%
%\ifthenelse{\equal{#2}{2}}{\hfill{ \ding{172}  \ding{183} \ding{174} } \addcontentsline{toc}{section}{#1 {\color{purple}\ding{183}}} }{%           
%\hfill{ \ding{172}  \ding{173} \ding{184} } \addcontentsline{toc}{section}{#1 {\color{orange}\ding{184}}}% 
%}%
%}%
%}
%}


%%%%%%%%%%%%%%%%%%%%% Cycle 4
\newcommand{\Titre}[2]{\section*{#1 
\ifthenelse{\equal{#2}{1}}   {\hfill{ \ding{182}  \ding{173} \ding{174} } \addcontentsline{toc}{section}{#1 \, \ding{182}} }%
{% sinon
\ifthenelse{\equal{#2}{1,5}}   {\hfill{ \ding{182}  \ding{183} \ding{174} } \addcontentsline{toc}{section}{#1 \, \ding{182} {\color{purple}\ding{183}}} }%
{% sinon
\ifthenelse{\equal{#2}{2}}   {\hfill{ \ding{172}  \ding{183} \ding{174} } \addcontentsline{toc}{section}{#1 \, {\color{purple}\ding{183}}} }
{% sinon
\ifthenelse{\equal{#2}{2,5}}   {\hfill{ \ding{172}  \ding{183} \ding{184} } \addcontentsline{toc}{section}{#1 \, {\color{purple}\ding{183}}  {\color{orange}\ding{184}}} }%
{% sinon
\hfill{ \ding{172}  \ding{173} \ding{184} } \addcontentsline{toc}{section}{#1 \,{\color{orange}\ding{184}}}% 
}%
}%
}%
}%
}%
}

%%%%%%%%%%%%% Titre
\newenvironment{titre}[2][]{%
\vspace{0.5cm}
\begin{tcolorbox}[enhanced, lifted shadow={0mm}{0mm}{0mm}{0mm}%
{black!60!white}, attach boxed title to top left={xshift=110mm, yshift*=-3mm}, coltitle=violet, colback=bleu3!25!white, boxed title style={colback=white!100}, colframe=bleu3,title=\stepcounter{cpttitre} \textbf{Fiche \thecpttitre}. #1 #2 ]}
{%
\end{tcolorbox}
\par}



%%%%%%%%%%%%% Définitions
\newenvironment{Def}[1][]{%
\medskip \begin{tcolorbox}[widget,colback=sacado_violet!0,colframe=sacado_violet!75,
adjusted title= \stepcounter{cptdef} Définition \thecptdef . {#1} ]}
{%
\end{tcolorbox}\par}


\newenvironment{DefT}[2][]{%
\medskip \begin{tcolorbox}[widget,colback=sacado_violet!0,colframe=sacado_violet!75,
adjusted title= \stepcounter{cptdef} Définition \thecptdef . {#1} \textit{#2}]}
{%
\end{tcolorbox}\par}

%%%%%%%%%%%%% Proposition
\newenvironment{Prop}[1][]{%
\medskip \begin{tcolorbox}[widget,colback=sacado_blue!0,colframe=sacado_blue!75!black,
adjusted title= \stepcounter{cpttheo} Proposition \thecpttheo . {#1} ]}
{%
\end{tcolorbox}\par}

%%%%%%%%%%%%% Propriétés
\newenvironment{Pp}[1][]{%
\medskip \begin{tcolorbox}[widget,colback=sacado_blue!0,colframe=sacado_blue!75!black,
adjusted title= \stepcounter{cpttheo} Propriété \thecpttheo . {#1}]}
{%
\end{tcolorbox}\par}

\newenvironment{PpT}[2][]{%
\medskip \begin{tcolorbox}[widget,colback=sacado_blue!0,colframe=sacado_blue!75!black,
adjusted title= \stepcounter{cpttheo} Propriété \thecpttheo . {#1} #2]}
{%
\end{tcolorbox}\par}

\newenvironment{Pps}[1][]{%
\medskip \begin{tcolorbox}[widget,colback=sacado_blue!0,colframe=sacado_blue!75!black,
adjusted title= \stepcounter{cpttheo} Propriétés \thecpttheo . {#1}]}
{%
\end{tcolorbox}\par}

%%%%%%%%%%%%% Théorèmes
\newenvironment{ThT}[2][]{% théorème avec titre
\medskip \begin{tcolorbox}[widget,colback=sacado_blue!0,colframe=sacado_blue!75!black,
adjusted title= \stepcounter{cpttheo} Théorème \thecpttheo . {#1} #2]}
{%
\end{tcolorbox}\par}

\newenvironment{Th}[1][]{%
\medskip \begin{tcolorbox}[widget,colback=sacado_blue!0,colframe=sacado_blue!75!black,
adjusted title= \stepcounter{cpttheo} Théorème \thecpttheo . {#1}]}
{%
\end{tcolorbox}\par}

%%%%%%%%%%%%% Règles
\newenvironment{Reg}[1][]{%
\medskip \begin{tcolorbox}[widget,colback=sacado_blue!0,colframe=sacado_blue!75!black,
adjusted title= \stepcounter{cpttheo} Règle \thecpttheo . {#1}]}
{%
\end{tcolorbox}\par}

%%%%%%%%%%%%% REMARQUES
\newenvironment{Rq}[1][]{%
\begin{bclogo}[couleur=sacado_orange!0, arrondi =0.15, noborder=true, couleurBarre=sacado_orange, logo = \bcinfo ]{ 
{\color{info}\normalsize{Remarque#1}}}}
{%
\end{bclogo}
\par}


\newenvironment{Rqs}[1][]{%
\begin{bclogo}[couleur=sacado_orange!0, arrondi =0.15, noborder=true, couleurBarre=sacado_orange, logo = \bcinfo ]{ 
{\color{info}\normalsize{Remarques#1}}}}
{%
\end{bclogo}
\par}

%%%%%%%%%%%%% EXEMPLES
\newenvironment{Ex}[1][]{%
\begin{bclogo}[couleur=sacado_yellow!15, arrondi =0.15, noborder=true, couleurBarre=sacado_yellow, logo = \bclampe ]{ 
\normalsize{Exemple#1}}}
{%
\end{bclogo}
\par}




%%%%%%%%%%%%% Preuve
\newenvironment{Pv}[1][]{%
\begin{tcolorbox}[breakable, enhanced,widget, colback=sacado_blue!10!white,boxrule=0pt,frame hidden,
borderline west={1mm}{0mm}{sacado_blue!75}]
\textbf{Preuve#1 : }}
{%
\end{tcolorbox}
\par}


%%%%%%%%%%%%% PreuveROC
\newenvironment{PvR}[1][]{%
\begin{tcolorbox}[breakable, enhanced,widget, colback=sacado_blue!10!white,boxrule=0pt,frame hidden,
borderline west={1mm}{0mm}{sacado_blue!75}]
\textbf{Preuve (ROC)#1 : }}
{%
\end{tcolorbox}
\par}


%%%%%%%%%%%%% Compétences
\newenvironment{Cps}[1][]{%
\vspace{0.4cm}
\begin{tcolorbox}[enhanced, lifted shadow={0mm}{0mm}{0mm}{0mm}%
{black!60!white}, attach boxed title to top left={xshift=5mm, yshift*=-3mm}, coltitle=white, colback=white, boxed title style={colback=sacado_green!100}, colframe=sacado_green!75!black,title=\textbf{Compétences associées#1}]}
{%
\end{tcolorbox}
\par}

%%%%%%%%%%%%% Compétences Collège
\newenvironment{CpsCol}[1][]{%
\vspace{0.4cm}
\begin{tcolorbox}[breakable, enhanced,widget, colback=white ,boxrule=0pt,frame hidden,
borderline west={2mm}{0mm}{bleu3}]
\textbf{#1}}
{%
\end{tcolorbox}
\par}




%%%%%%%%%%%%% Attendus
\newenvironment{Ats}[1][]{%
\vspace{0.4cm}
\begin{tcolorbox}[enhanced, lifted shadow={0mm}{0mm}{0mm}{0mm}%
{black!60!white}, attach boxed title to top left={xshift=5mm, yshift*=-3mm}, coltitle=white, colback=white, boxed title style={colback=sacado_green!100}, colframe=sacado_green!75!black,title=\textbf{Attendus du chapitre#1}]}
{%
\end{tcolorbox}
\par}

%%%%%%%%%%%%% Méthode
\newenvironment{Mt}[1][]{%
\vspace{0.4cm}
\begin{bclogo}[couleur=sacado_blue!0, arrondi =0.15, noborder=true, couleurBarre=bleu3, logo = \bccrayon ]{ 
\normalsize{{\color{bleu3}Méthode #1}}}}
{%
\end{bclogo}
\par}


%%%%%%%%%%%%% Méthode en vidéo
\newcommand{\MtV}[2]{\vspace{0.4cm} \colorbox{sacado_blue!0}{\hspace{0.2 cm}\tikz\node[rounded corners=1pt,draw] {\color{red}$\blacktriangleright$}; \quad  \href{https://youtu.be/#1?rel=0}{\raisebox{0.8mm}{{\color{red}\textbf{Méthode en vidéo : #2}}}}}}


%%%%%%%%%%%%% A voir (AV) : Lien externe + vidéo non Youtube
\newcommand{\AV}[2]{\vspace{0.4cm} \colorbox{bleu1!0}{\hspace{0.2 cm}\tikz\node[rounded corners=1pt,draw] {\color{red}$\blacktriangleright$}; \quad  \href{#1}{\raisebox{0.8mm}{{\color{red}\textbf{#2}}}}}}


%%%%%%%%%%%%% Etymologie
\newenvironment{Ety}[1][]{%
\begin{bclogo}[couleur=sacado_green!0, arrondi =0.15, noborder=true, couleurBarre=sacado_green, logo = \bcplume ]{ 
\normalsize{{\color{sacado_green}Étymologie#1}}}}
{%
\end{bclogo}
\par}


%%%%%%%%%%%%% Notation
\newenvironment{Nt}[1][]{%
\begin{bclogo}[couleur=sacado_violet!0, arrondi =0.15, noborder=true, couleurBarre=sacado_violet!75, logo = \bccrayon ]{ 
\normalsize{{\color{violet!75}Notation#1}}}}
{%
\end{bclogo}
\par}
%%%%%%%%%%%%% Histoire
%\newenvironment{His}[1][]{%
%\begin{bclogo}[couleur=brown!30, arrondi =0.15, noborder=true, couleurBarre=brown, logo = \bcvaletcoeur ]{ 
%\normalsize{{\color{brown}Histoire des mathématiques#1}}}}
%{%
%\end{bclogo}
%\par}

\newenvironment{His}[1][]{%
\vspace{0.4cm}
\begin{tcolorbox}[enhanced, lifted shadow={0mm}{0mm}{0mm}{0mm}%
{brown!60!white}, attach boxed title to top left={xshift=5mm, yshift*=-3mm}, coltitle=white, colback=white, boxed title style={colback=brown!100}, colframe=brown!75!black,title=\textbf{Histoire des mathématiques#1}]}
{%
\end{tcolorbox}
\par}

%%%%%%%%%%%%% Attention
\newenvironment{Att}[1][]{%
\begin{bclogo}[couleur=red!0, arrondi =0.15, noborder=true, couleurBarre=red, logo = \bcattention ]{ 
\normalsize{{\color{red}Attention. #1}}}}
{%
\end{bclogo}
\par}


%%%%%%%%%%%%% Conséquence
\newenvironment{Cq}[1][]{%
\textbf{Conséquence #1}}
{%
\par}

%%%%%%%%%%%%% Vocabulaire
\newenvironment{Voc}[1][]{%
\setlength{\logowidth}{10pt}
%\begin{footnotesize}
\begin{bclogo}[ noborder , couleur=white, logo =\bcbook]{#1}}
{%
\end{bclogo}
%\end{footnotesize}
\par}


%%%%%%%%%%%%% Video
\newenvironment{Vid}[1][]{%
\setlength{\logowidth}{12pt}
\begin{bclogo}[ noborder , couleur=white,barre=none, logo =\bcoeil]{#1}}
{%
\end{bclogo}
\par}


%%%%%%%%%%%%% Syntaxe
\newenvironment{Syn}[1][]{%
\begin{bclogo}[couleur=violet!0, arrondi =0.15, noborder=true, couleurBarre=violet!75, logo = \bcicosaedre ]{ 
\normalsize{{\color{violet!75}Syntaxe#1}}}}
{%
\end{bclogo}
\par}

%%%%%%%%%%%%% Auto évaluation
\newenvironment{autoeval}[1][]{%
\vspace{0.4cm}
\begin{tcolorbox}[enhanced, lifted shadow={0mm}{0mm}{0mm}{0mm}%
{black!60!white}, attach boxed title to top left={xshift=5mm, yshift*=-3mm}, coltitle=white, colback=white, boxed title style={colback=sacado_green!100}, colframe=sacado_green!75!black,title=\textbf{J'évalue mes compétences#1}]}
{%
\end{tcolorbox}
\par}


\newenvironment{autotest}[1][]{%
\vspace{0.4cm}
\begin{tcolorbox}[enhanced, lifted shadow={0mm}{0mm}{0mm}{0mm}%
{red!60!white}, attach boxed title to top left={xshift=5mm, yshift*=-3mm}, coltitle=white, colback=white, boxed title style={colback=red!100}, colframe=red!75!black,title=\textbf{Pour faire le point #1}]}
{%
\end{tcolorbox}
\par}

\newenvironment{ExOApp}[1][]{% Exercice d'application direct
\vspace{0.4cm}
\begin{tcolorbox}[enhanced, lifted shadow={0mm}{0mm}{0mm}{0mm}%
{red!60!white}, attach boxed title to top left={xshift=5mm, yshift*=-3mm}, coltitle=white, colback=white, boxed title style={colback=sacado_green!100}, colframe=sacado_green!75!black,title=\textbf{Application #1}]}
{%
\end{tcolorbox}
\par}

\newenvironment{ExOInt}[1][]{% Exercice d'application direct
\vspace{0.4cm}
\begin{tcolorbox}[enhanced, lifted shadow={0mm}{0mm}{0mm}{0mm}%
{red!60!white}, attach boxed title to top left={xshift=5mm, yshift*=-3mm}, coltitle=white, colback=white, boxed title style={colback=sacado_green!50}, colframe=sacado_green!75!black,title=\textbf{Exercice #1}]}
{%
\end{tcolorbox}
\par}

%Illustrations
\newtcolorbox{Illqr}[1]{
  enhanced,
  colback=white,
  colframe=ill_frame,
  colbacktitle=ill_back,
  coltitle=ill_title,
  title=\textbf{Illustration},
  boxrule=1pt, % épaisseur du trait à 1pt
  center,
  overlay={
    \node[anchor=south east, inner sep=0pt,xshift=-1pt,yshift=2pt,fill=white] at (frame.south east) {\fancyqr[height=1cm]{#1}};
  },
  after=\par,
  before=\vspace{0.4cm},
}

\newtcolorbox{Ill}{
  enhanced,
  colback=white,
  colframe=ill_frame,
  colbacktitle=ill_back,
  coltitle=ill_title,
  title=\textbf{Illustration},
  boxrule=1pt, % épaisseur du trait à 1pt
  center,
  after=\par,
  before=\vspace{0.4cm},
}

%%%%%%%%%%%%%% Propriétés
%\newenvironment{Pp}[1][]{%
%\medskip \begin{tcolorbox}[widget,colback=sacado_blue!0,colframe=sacado_blue!75!black,
%adjusted title= \stepcounter{cpttheo} Propriété \thecpttheo . {#1}]}
%{%
%\end{tcolorbox}\par}

%%%%% Pour réinitialiser numéros des chapitres après une nouvelle partie
% \makeatletter
    % \@addtoreset{section}{part}
% \makeatother

% \newcommand{\EPC}[3]{ % Exercice par compétence de niveau 1
% \ifthenelse{\equal{#1}{1}}
% {%condition2 vraie
% \vspace{0.4cm}
% \stepcounter{cptex}
% \tikz\node[rounded corners=0pt,draw,fill=bleu2]{\color{white}\textbf{ \thecptex}}; \quad  {\color{bleu2}\textbf{#3}}
% \input{#2}
% }% fin condition2 vraie
% {%condition2 fausse
% \vspace{0.4cm}
% \stepcounter{cptex}
% \tikz\node[rounded corners=2pt,draw,fill=eduscol4P]{\color{white}\textbf{ \thecptex}}; \quad  {\color{eduscol4P} \textbf{En temps libre.} \textbf{ #3}} 
% \input{#2}
% }% fin condition2 fausse
% } % fin de la procédure
%
%\usepackage{hyperref}
%
%\begin{document}
%
%%-------------------------------
%%	TITLE SECTION
%%-------------------------------
%
%\fancyhead[C]{}
%\hrule\medskip % Upper rule
%\begin{minipage}{0.295\textwidth} 
%\raggedright
%Classe \myClasse \hfill\\
%\myDiscipline \hfill\\
%\myParcours \hfill\\
%\end{minipage}
%\begin{minipage}{0.4\textwidth} 
%\centering 
%\scshape\huge
%\textcolor{sacado_purple}{\myTitle} \\ 
%\normalsize 
%%\mySubTitle \\ 
%\end{minipage}
%\begin{minipage}{0.295\textwidth} 
%\raggedleft
%\href{https://sacado.xyz/}{\includegraphics[width=.2\linewidth]{sacadoA1.png}}
%%\myAnnee \hfill\\
%\end{minipage}
%\medskip \hrule
%\bigskip

%-------------------------------
%	CONTENTS
%-------------------------------

\chapter{Polygones et cercles}
{https://sacado.xyz/qcm/parcours_show_course/0/117131}

\begin{pageCours}

\section{Polygones}

\mini{.48\linewidth}{\begin{Def}
Un \textbf{polygone} est une \textbf{figure} géométrique \textbf{fermée} dont les \textbf{côtés} sont des \textbf{segments} de droites.
\end{Def}}
{.48\linewidth}{
\begin{Ill}
Le polygone $ABCDE$ a $5$ sommets et $5$ côtés.

\begin{tikzpicture}[line cap=round,line join=round,>=triangle 45,x=.8cm,y=.8cm]
\clip(-2,-2) rectangle (5,4);
\fill[line width=2.pt,color=zzttqq,fill=zzttqq,fill opacity=0.10000000149011612] (-1.063311344808564,1.6089901265883788) -- (1.4923210474495676,2.7760130109143466) -- (2.1127635935469176,0.9442302557697901) -- (4.2695400633138965,0.) -- (2.467302191316832,-1.330725746587158) -- cycle;
\draw [line width=2.pt,color=zzttqq] (-1.063311344808564,1.6089901265883788)-- (1.4923210474495676,2.7760130109143466);
\draw [line width=2.pt,color=zzttqq] (1.4923210474495676,2.7760130109143466)-- (2.1127635935469176,0.9442302557697901);
\draw [line width=2.pt,color=zzttqq] (2.1127635935469176,0.9442302557697901)-- (4.2695400633138965,0.);
\draw [line width=2.pt,color=zzttqq] (4.2695400633138965,0.)-- (2.467302191316832,-1.330725746587158);
\draw [line width=2.pt,color=zzttqq] (2.467302191316832,-1.330725746587158)-- (-1.063311344808564,1.6089901265883788);
\draw[color=ududff] (-0.9599042537923393,1.8822802957026874) node {$A$};
\draw[color=ududff] (1.5957281384657929,3.049303180028654) node {$B$};
\draw[color=ududff] (2.2161706845631426,1.2175204248840987) node {$C$};
\draw[color=xdxdff] (4.372947154330122,0.272084164164328) node {$D$};
\draw[color=ududff] (2.570709282333057,-1.0574355774728494) node {$E$};
\end{tikzpicture}
\end{Ill}}

\mini{.58\linewidth}{\begin{DefT}{Polygones réguliers}
Un polygone est dit \textbf{régulier} lorsque tous ses côtés sont de \textbf{même longueur} et tous ses \textbf{angles} ont la \textbf{même} mesure.
\end{DefT}
\begin{Rq}
Il existe un \textbf{cercle} qui passe par \textbf{tous les sommets} d'un polygone régulier.
\end{Rq}}
{.38\linewidth}{\begin{Illqr}{Appliquette poly régulier}
\begin{center}
\begin{tikzpicture}[line cap=round,line join=round,>=triangle 45,x=0.8cm,y=0.8cm]
\clip(-3,-3) rectangle (3,3);
\fill[line width=2.pt,color=zzttqq,fill=zzttqq,fill opacity=0.10000000149011612] (2.,0.) -- (-1.,1.7320508075688772) -- (-1.,-1.7320508075688774) -- cycle;
\draw [shift={(2.,0.)},line width=2.pt,color=xdxdff,fill=xdxdff,fill opacity=0.10000000149011612] (0,0) -- (150.:0.5323591126142355) arc (150.:210.:0.5323591126142355) -- cycle;
\draw [shift={(-1.,-1.7320508075688774)},line width=2.pt,color=xdxdff,fill=xdxdff,fill opacity=0.10000000149011612] (0,0) -- (30.:0.5323591126142355) arc (30.:90.:0.5323591126142355) -- cycle;
\draw [shift={(-1.,1.7320508075688772)},line width=2.pt,color=xdxdff,fill=xdxdff,fill opacity=0.10000000149011612] (0,0) -- (-90.:0.5323591126142355) arc (-90.:-30.:0.5323591126142355) -- cycle;
\draw [line width=2.pt] (0.,0.) circle (1.6cm);
\draw [line width=2.pt,color=zzttqq] (2.,0.)-- (-1.,1.7320508075688772);
\draw [line width=2.pt,color=zzttqq] (-1.,1.7320508075688772)-- (-1.,-1.7320508075688774);
\draw [line width=2.pt,color=zzttqq] (-1.,-1.7320508075688774)-- (2.,0.);
\draw [line width=1.pt,dash pattern=on 1pt off 3pt] (2.,0.)-- (0.,0.);
\draw [line width=1.pt,dash pattern=on 1pt off 3pt] (-1.,-1.7320508075688774)-- (0.,0.);
\draw [line width=1.pt,dash pattern=on 1pt off 3pt] (-1.,1.7320508075688772)-- (0.,0.);
\draw[color=xdxdff] (2.073110611619774,0.19681571921764635) node {$A$};
\draw[color=ududff] (-0.9293947835245144,1.9323064263400531) node {$B$};
\draw[color=qqqqff] (-0.9293947835245144,-1.5386749879047603) node {$C$};
\end{tikzpicture}
\end{center}
\end{Illqr}}

\section{Triangles}

\mini{.38\linewidth}{\begin{Def}
Un \textbf{triangle} est un polygone ayant \textbf{3 côtés}.

Un triangle a 3 sommets.
\end{Def}}
{.68\linewidth}{
\begin{Ill}
\centering
\begin{tikzpicture}[line cap=round,line join=round,>=triangle 45,x=0.8cm,y=0.8cm]
\clip(-3,-3) rectangle (3,3);
\fill[line width=2.pt,color=zzttqq,fill=zzttqq,fill opacity=0.10000000149011612] (-2.270939747312388,-0.20245361524303004) -- (-0.23732793712600822,2.12927929800732) -- (2.,-2.) -- cycle;
\draw [shift={(2.,-2.)},line width=2.pt,color=violet,fill=violet,fill opacity=0.10000000149011612] (0,0) -- (118.44974022885543:0.47912320135281194) arc (118.44974022885543:157.1748246920994:0.47912320135281194) -- cycle;
\draw [shift={(-2.270939747312388,-0.20245361524303004)},line width=2.pt,color=xdxdff,fill=xdxdff,fill opacity=0.10000000149011612] (0,0) -- (-22.82517530790062:0.47912320135281194) arc (-22.82517530790062:48.9068163841034:0.47912320135281194) -- cycle;
\draw [shift={(-0.23732793712600822,2.12927929800732)},line width=2.pt,color=qqwuqq,fill=qqwuqq,fill opacity=0.10000000149011612] (0,0) -- (-131.09318361589658:0.47912320135281194) arc (-131.09318361589658:-61.55025977114456:0.47912320135281194) -- cycle;
\draw [line width=2.pt,color=zzttqq] (-2.270939747312388,-0.20245361524303004)-- (-0.23732793712600822,2.12927929800732);
\draw [line width=2.pt,color=zzttqq] (-0.23732793712600822,2.12927929800732)-- (2.,-2.);
\draw [line width=2.pt,color=zzttqq] (2.,-2.)-- (-2.270939747312388,-0.20245361524303004);
\draw[color=ududff] (-2.196409471546395,-0.005480743575763017) node {$A$};
\draw[color=ududff] (-0.16279766136001528,2.326252169674587) node {$B$};
\draw[color=ududff] (2.073110611619774,-1.8048545442118777) node {$C$};
\end{tikzpicture}
\end{Ill}}

\subsection{Triangles particuliers}

\mini{.48\linewidth}{
\begin{DefT}{Triangle équilatéral}
Un triangle \textbf{équilatéral} est un triangle dont les trois côtés ont la même mesure.
\end{DefT}
\begin{Pp}
Dans un triangle équilatéral les \textbf{trois} angles ont la même mesure : 60°.
\end{Pp}
}{.48\linewidth}{
\begin{Ill}
\centering
\begin{tikzpicture}[line cap=round,line join=round,>=triangle 45,x=0.8cm,y=0.8cm]
\clip(-3,-3) rectangle (3,3);
\fill[line width=2.pt,color=zzttqq,fill=zzttqq,fill opacity=0.10000000149011612] (-1.3552820736159028,-1.5972344902923261) -- (2.04116906486292,-0.5112219005592864) -- (-0.5975709959150303,1.8871847732094331) -- cycle;
\draw [shift={(2.04116906486292,-0.5112219005592864)},line width=2.pt,color=xdxdff,fill=xdxdff,fill opacity=0.10000000149011612] (0,0) -- (137.73163652585896:0.5323591126142355) arc (137.73163652585896:197.731636525859:0.5323591126142355) -- cycle;
\draw [shift={(-1.3552820736159028,-1.5972344902923261)},line width=2.pt,color=xdxdff,fill=xdxdff,fill opacity=0.10000000149011612] (0,0) -- (17.731636525858974:0.5323591126142355) arc (17.731636525858974:77.73163652585897:0.5323591126142355) -- cycle;
\draw [shift={(-0.5975709959150303,1.8871847732094331)},line width=2.pt,color=xdxdff,fill=xdxdff,fill opacity=0.10000000149011612] (0,0) -- (-102.26836347414104:0.5323591126142355) arc (-102.26836347414104:-42.26836347414104:0.5323591126142355) -- cycle;
\draw [line width=2.pt,color=zzttqq] (-1.3552820736159028,-1.5972344902923261)-- (2.04116906486292,-0.5112219005592864);
\draw [line width=2.pt,color=zzttqq] (2.04116906486292,-0.5112219005592864)-- (-0.5975709959150303,1.8871847732094331);
\draw [line width=2.pt,color=zzttqq] (-0.5975709959150303,1.8871847732094331)-- (-1.3552820736159028,-1.5972344902923261);
\draw[color=ududff] (-1.674697541184444,-1.3896144363727745) node {$A$};
\draw[color=ududff] (2.115699340628913,-0.31424902889201944) node {$B$};
\draw[color=qqqqff] (-0.5248018579376954,2.081366977872039) node {$C$};
\draw[color=xdxdff] (1.1787473024278583,-0.38877930465801236) node {$60\textrm{\degre}$};
\end{tikzpicture}
\end{Ill}
}

\mini{.48\linewidth}{
\begin{DefT}{Triangle isocèle}
Un triangle \textbf{isocèle} est un triangle dont deux côtés ont la même mesure.
\end{DefT}
\begin{Pp}
Un triangle isocèle posséde \textbf{deux} angles de \textbf{même mesure}.
\end{Pp}
}{.48\linewidth}{
\begin{Ill}
\centering
\begin{tikzpicture}[line cap=round,line join=round,>=triangle 45,x=0.8cm,y=0.8cm]
\clip(-3,-3) rectangle (3,3);
\fill[line width=2.pt,color=zzttqq,fill=zzttqq,fill opacity=0.10000000149011612] (-1.5256369896524582,-1.171347200200938) -- (1.9027556955832186,-2.35318443020454) -- (1.479155321328516,1.9816251921209547) -- cycle;
\draw [shift={(1.9027556955832186,-2.35318443020454)},line width=2.pt,color=xdxdff,fill=xdxdff,fill opacity=0.10000000149011612] (0,0) -- (95.58125955604939:0.31941546756854133) arc (95.58125955604939:160.97987786167886:0.31941546756854133) -- cycle;
\draw [shift={(-1.5256369896524582,-1.171347200200938)},line width=2.pt,color=xdxdff,fill=xdxdff,fill opacity=0.10000000149011612] (0,0) -- (-19.02012213832115:0.31941546756854133) arc (-19.02012213832115:46.378496167308306:0.31941546756854133) -- cycle;
\draw [shift={(1.479155321328516,1.9816251921209547)},line width=2.pt,color=qqwuqq,fill=qqwuqq,fill opacity=0.10000000149011612] (0,0) -- (-133.6215038326917:0.31941546756854133) arc (-133.6215038326917:-84.41874044395061:0.31941546756854133) -- cycle;
\draw [line width=2.pt,color=zzttqq] (-1.5256369896524582,-1.171347200200938)-- (1.9027556955832186,-2.35318443020454);
\draw [line width=2.pt,color=zzttqq] (1.9027556955832186,-2.35318443020454)-- (1.479155321328516,1.9816251921209547);
\draw [line width=2.pt,color=zzttqq] (1.479155321328516,1.9816251921209547)-- (-1.5256369896524582,-1.171347200200938);
\draw[color=ududff] (-1.8450524572209994,-0.9637271462813863) node {$A$};
\draw[color=ududff] (1.9772859713492117,-2.156211558537273) node {$B$};
\draw[color=xdxdff] (1.551398681257823,2.177191618142601) node {$C$};
\end{tikzpicture}
\end{Ill}
}

\mini{.48\linewidth}{
\begin{DefT}{Triangle rectangle}
Un triangle \textbf{rectangle} est un triangle dont un \textbf{angle est droit} (sa mesure est de 90°).
\end{DefT}
\begin{Def}
Dans un triangle  rectangle, le côté opposé à l'angle droite s'appelle l' \textbf{hypoténuse} du triangle.
\end{Def}
}{.48\linewidth}{
\begin{Ill}
\centering
\begin{tikzpicture}[line cap=round,line join=round,>=triangle 45,x=0.8cm,y=0.8cm]
\clip(-3,-3) rectangle (3,3);
\fill[line width=2.pt,color=zzttqq,fill=zzttqq,fill opacity=0.10000000149011612] (-1.5256369896524582,-1.171347200200938) -- (2.2328183454040444,-0.6496352698389876) -- (1.8498215704497243,2.109504761974792) -- cycle;
\draw[line width=2.pt,color=xdxdff,fill=xdxdff,fill opacity=0.10000000149011612] (2.1862372898057174,-0.3140615427734898) -- (1.85066356274022,-0.3606425983718165) -- (1.8972446183385467,-0.6962163254373143) -- (2.2328183454040444,-0.6496352698389876) -- cycle; 
\draw [line width=2.pt,color=zzttqq] (-1.5256369896524582,-1.171347200200938)-- (2.2328183454040444,-0.6496352698389876);
\draw [line width=2.pt,color=zzttqq] (2.2328183454040444,-0.6496352698389876)-- (1.8498215704497243,2.109504761974792);
\draw [line width=2.pt,color=zzttqq] (1.8498215704497243,2.109504761974792)-- (-1.5256369896524582,-1.171347200200938);
\draw[color=ududff] (-1.8450524572209994,-0.9637271462813863) node {$A$};
\draw[color=ududff] (2.307348621170038,-0.45266239817172055) node {$B$};
\draw[color=xdxdff] (1.9240500600877881,2.3049578051700172) node {$C$};
\draw[color=xdxdff] (1.6046345925192467,-0.2610131176305959) node {$90\textrm{\degre}$};
\end{tikzpicture}
\end{Ill}
}

\subsection{Construire un triangle}

\begin{Mt}
\qr{•}{Construire un triangle dont on connait la longueur des trois côtés}
\end{Mt}

\section{Quadrilatères}

\mini{.38\linewidth}{\begin{Def}
Un quadrilatère est un polygone qui a 4 côtés.

$ABCD$ est un quadrilatère.

Les segments $[AC]$ et $[BD]$ sont les diagonales de $ABCD$.
\end{Def}}
{.68\linewidth}{
\begin{Illqr}{Quadrilatères avec une particularité}
\centering
\begin{tikzpicture}[line cap=round,line join=round,>=triangle 45,x=0.8cm,y=0.8cm]
\clip(-3,-3) rectangle (3,3);
\fill[line width=2.pt,color=zzttqq,fill=zzttqq,fill opacity=0.10000000149011612] (-1.5575785364093122,-1.171347200200938) -- (-1.0997496995610698,1.4265652693565296) -- (2.3818788969360307,1.8205110126910635) -- (2.296701438917753,-0.22374797974759944) -- cycle;
\draw [line width=2.pt,color=zzttqq] (-1.5575785364093122,-1.171347200200938)-- (-1.0997496995610698,1.4265652693565296);
\draw [line width=2.pt,color=zzttqq] (-1.0997496995610698,1.4265652693565296)-- (2.3818788969360307,1.8205110126910635);
\draw [line width=2.pt,color=zzttqq] (2.3818788969360307,1.8205110126910635)-- (2.296701438917753,-0.22374797974759944);
\draw [line width=2.pt,color=zzttqq] (2.296701438917753,-0.22374797974759944)-- (-1.5575785364093122,-1.171347200200938);
\draw [line width=1.pt,dash pattern=on 1pt off 3pt] (-1.5575785364093122,-1.171347200200938)-- (2.3818788969360307,1.8205110126910635);
\draw [line width=1.pt,dash pattern=on 1pt off 3pt] (-1.0997496995610698,1.4265652693565296)-- (2.296701438917753,-0.22374797974759944);
\draw[color=ududff] (-1.4830482606433193,-0.974374328533671) node {$A$};
\draw[color=ududff] (-1.0252194237950767,1.6235381410237966) node {$B$};
\draw[color=ududff] (2.4564091727020236,2.0174838843583305) node {$C$};
\draw[color=ududff] (2.371231714683746,-0.026775108080332427) node {$D$};
\end{tikzpicture}
\end{Illqr}
\begin{flushright}
Scan le QR code pour les particularités qu'on peut trouver sur les quadrilatères.
\end{flushright}
}

\subsection{Quadrilatères particuliers}

\mini{.48\linewidth}{
\begin{DefT}{Losange}
Un \textbf{losange} est un quadrilatère qui a ses \textbf{quatre côtés de même longueur}.
\end{DefT}
\begin{Pp}
Les \textbf{diagonales} d'un losange sont \textbf{perpendiculaires}.
\end{Pp}
}{.48\linewidth}{
\begin{Ill}
\centering
\begin{tikzpicture}[line cap=round,line join=round,>=triangle 45,x=0.8cm,y=0.8cm]
\clip(-3,-3) rectangle (3,3);
\fill[line width=2.pt,color=zzttqq,fill=zzttqq,fill opacity=0.10000000149011612] (-1.,1.) -- (2.,2.) -- (1.,-1.) -- (-2.,-2.) -- cycle;
\draw[line width=2.pt,color=xdxdff,fill=xdxdff,fill opacity=0.10000000149011612] (0.15970773378427064,0.1597077337842706) -- (0.,0.3194154675685413) -- (-0.1597077337842706,0.15970773378427064) -- (0.,0.) -- cycle; 
\draw [line width=1.pt,dash pattern=on 1pt off 3pt] (-1.,1.)-- (1.,-1.);
\draw [line width=1.pt,dash pattern=on 1pt off 3pt] (-2.,-2.)-- (2.,2.);
\draw [line width=2.pt,color=zzttqq] (-1.,1.)-- (2.,2.);
\draw [line width=2.pt,color=zzttqq] (0.45454638211680987,1.5521874871992187) -- (0.4949495980129798,1.430977839510711);
\draw [line width=2.pt,color=zzttqq] (0.5050504019870204,1.569022160489289) -- (0.5454536178831904,1.4478125128007813);
\draw [line width=2.pt,color=zzttqq] (2.,2.)-- (1.,-1.);
\draw [line width=2.pt,color=zzttqq] (1.5690221604892896,0.5050504019870207) -- (1.4478125128007822,0.5454536178831906);
\draw [line width=2.pt,color=zzttqq] (1.5521874871992185,0.454546382116809) -- (1.4309778395107111,0.49494959801297883);
\draw [line width=2.pt,color=zzttqq] (1.,-1.)-- (-2.,-2.);
\draw [line width=2.pt,color=zzttqq] (-0.45454638211680987,-1.5521874871992187) -- (-0.4949495980129798,-1.4309778395107113);
\draw [line width=2.pt,color=zzttqq] (-0.5050504019870204,-1.5690221604892898) -- (-0.5454536178831904,-1.4478125128007824);
\draw [line width=2.pt,color=zzttqq] (-2.,-2.)-- (-1.,1.);
\draw [line width=2.pt,color=zzttqq] (-1.5690221604892896,-0.5050504019870207) -- (-1.4478125128007822,-0.5454536178831906);
\draw [line width=2.pt,color=zzttqq] (-1.5521874871992185,-0.454546382116809) -- (-1.4309778395107111,-0.49494959801297883);
\draw[color=ududff] (-0.9293947835245144,1.1976508509324084) node {$A$};
\draw[color=ududff] (1.0722754799050112,-0.8040194124971158) node {$B$};
\draw[color=ududff] (-1.930229915239277,-1.8048545442118777) node {$C$};
\draw[color=ududff] (2.073110611619774,2.1984859826471705) node {$D$};
\end{tikzpicture}
\end{Ill}
}

\mini{.48\linewidth}{
\begin{DefT}{Rectangle}
Un \textbf{rectangle} est un quadrilatère qui possède \textbf{quatre angles droits}.
\end{DefT}
\begin{Pp}
Les \textbf{diagonales} d'un rectangle ont la \textbf{même longueur}.
\end{Pp}
}{.48\linewidth}{
\begin{Ill}
\centering
\begin{tikzpicture}[line cap=round,line join=round,>=triangle 45,x=0.8cm,y=0.8cm]
\clip(-3,-3) rectangle (3,3);
\fill[line width=2.pt,color=zzttqq,fill=zzttqq,fill opacity=0.10000000149011612] (-2.,1.) -- (2.,1.) -- (2.,-1.) -- (-2.,-1.) -- cycle;
\draw[line width=2.pt,color=xdxdff,fill=xdxdff,fill opacity=0.10000000149011612] (1.7741391568664127,1.) -- (1.7741391568664127,0.7741391568664128) -- (2.,0.7741391568664128) -- (2.,1.) -- cycle; 
\draw[line width=2.pt,color=xdxdff,fill=xdxdff,fill opacity=0.10000000149011612] (2.,-0.7741391568664128) -- (1.7741391568664127,-0.7741391568664127) -- (1.7741391568664127,-1.) -- (2.,-1.) -- cycle; 
\draw[line width=2.pt,color=xdxdff,fill=xdxdff,fill opacity=0.10000000149011612] (-1.7741391568664127,-1.) -- (-1.7741391568664127,-0.7741391568664128) -- (-2.,-0.7741391568664128) -- (-2.,-1.) -- cycle; 
\draw[line width=2.pt,color=xdxdff,fill=xdxdff,fill opacity=0.10000000149011612] (-2.,0.7741391568664128) -- (-1.7741391568664127,0.7741391568664128) -- (-1.7741391568664127,1.) -- (-2.,1.) -- cycle; 
\draw [line width=2.pt,color=zzttqq] (-2.,1.)-- (2.,1.);
\draw [line width=2.pt,color=zzttqq] (2.,1.)-- (2.,-1.);
\draw [line width=2.pt,color=zzttqq] (2.,-1.)-- (-2.,-1.);
\draw [line width=2.pt,color=zzttqq] (-2.,-1.)-- (-2.,1.);
\draw [line width=1.pt,dash pattern=on 1pt off 3pt] (-2.,1.)-- (2.,-1.);
\draw [line width=1.pt,dash pattern=on 1pt off 3pt] (-2.,-1.)-- (2.,1.);
\draw[color=ududff] (-1.930229915239277,1.1976508509324084) node {$A$};
\draw[color=ududff] (2.073110611619774,1.1976508509324084) node {$B$};
\draw[color=ududff] (2.073110611619774,-0.8040194124971158) node {$C$};
\draw[color=ududff] (-1.930229915239277,-0.8040194124971158) node {$D$};
\end{tikzpicture}
\end{Ill}
}

\mini{.48\linewidth}{
\begin{DefT}{Carré}
Un \textbf{carré} est un quadrilatère qui possède \textbf{quatre angles droits} et qui a ses \textbf{quatre côtés de même longueur}.
\end{DefT}
\begin{Pp}
Les \textbf{diagonales} d'un carré sont \textbf{perpendiculaires} et de \textbf{même longueur}.
\end{Pp}
}{.48\linewidth}{
\begin{Ill}
\centering
\begin{tikzpicture}[line cap=round,line join=round,>=triangle 45,x=0.8cm,y=0.8cm]
\clip(-3,-3) rectangle (3,3);
\fill[line width=2.pt,color=zzttqq,fill=zzttqq,fill opacity=0.10000000149011612] (-0.39703567091027886,-2.0444161448882836) -- (2.0092275181060657,-0.18115925073846062) -- (0.14597062395624327,2.2251039382778837) -- (-2.2602925650601016,0.3618470441280611) -- cycle;
\draw[line width=2.pt,color=xdxdff,fill=xdxdff,fill opacity=0.10000000149011612] (-0.09213762139290618,2.0407280845783213) -- (0.0922382323066562,1.8026198392291717) -- (0.33034647765580566,1.9869956929287342) -- (0.14597062395624327,2.2251039382778837) -- cycle; 
\draw[line width=2.pt,color=xdxdff,fill=xdxdff,fill opacity=0.10000000149011612] (1.8248516644065031,0.05694899461068881) -- (1.5867434190573537,-0.12742685908887366) -- (1.7711192727569163,-0.36553510443802306) -- (2.0092275181060657,-0.18115925073846062) -- cycle; 
\draw[line width=2.pt,color=xdxdff,fill=xdxdff,fill opacity=0.10000000149011612] (-0.15892742556112943,-1.860040291188721) -- (-0.34330327926069193,-1.6219320458395716) -- (-0.5814115246098414,-1.8063078995391342) -- (-0.39703567091027886,-2.0444161448882836) -- cycle; 
\draw[line width=2.pt,color=xdxdff,fill=xdxdff,fill opacity=0.10000000149011612] (-2.0759167113605392,0.1237387987789117) -- (-1.8378084660113896,0.30811465247847414) -- (-2.022184319710952,0.5462228978276236) -- (-2.2602925650601016,0.3618470441280611) -- cycle; 
\draw[line width=2.pt,color=xdxdff,fill=xdxdff,fill opacity=0.10000000149011612] (-0.09228730228532271,0.3517425966530293) -- (-0.35368600224355173,0.3849878178447245) -- (-0.3869312234352469,0.12358911788649543) -- (-0.12553252347701782,0.09034389669480025) -- cycle; 
\draw [line width=2.pt,color=zzttqq] (-0.39703567091027886,-2.0444161448882836)-- (2.0092275181060657,-0.18115925073846062);
\draw [line width=2.pt,color=zzttqq] (0.7459379042986362,-1.0785739883852803) -- (0.8241619541614476,-1.1795947613509685);
\draw [line width=2.pt,color=zzttqq] (0.7880298930343392,-1.0459806342757758) -- (0.8662539428971506,-1.147001407241464);
\draw [line width=2.pt,color=zzttqq] (2.0092275181060657,-0.18115925073846062)-- (0.14597062395624327,2.2251039382778837);
\draw [line width=2.pt,color=zzttqq] (1.0433853616030633,0.9618143244704542) -- (1.1444061345687517,1.0400383743332657);
\draw [line width=2.pt,color=zzttqq] (1.0107920074935577,1.0039063132061572) -- (1.1118127804592461,1.0821303630689687);
\draw [line width=2.pt,color=zzttqq] (0.14597062395624327,2.2251039382778837)-- (-2.2602925650601016,0.3618470441280611);
\draw [line width=2.pt,color=zzttqq] (-0.9970029512526712,1.2592617817748812) -- (-1.0752270011154828,1.3602825547405692);
\draw [line width=2.pt,color=zzttqq] (-1.0390949399883742,1.2266684276653763) -- (-1.1173189898511857,1.3276892006310645);
\draw [line width=2.pt,color=zzttqq] (-2.2602925650601016,0.3618470441280611)-- (-0.39703567091027886,-2.0444161448882836);
\draw [line width=2.pt,color=zzttqq] (-1.2944504085570996,-0.7811265310808538) -- (-1.395471181522788,-0.8593505809436652);
\draw [line width=2.pt,color=zzttqq] (-1.261857054447594,-0.8232185198165568) -- (-1.3628778274132824,-0.9014425696793682);
\draw [line width=1.pt,dash pattern=on 1pt off 3pt] (-2.2602925650601016,0.3618470441280611)-- (2.0092275181060657,-0.18115925073846062);
\draw [line width=1.pt,dash pattern=on 1pt off 3pt] (0.14597062395624327,2.2251039382778837)-- (-0.39703567091027886,-2.0444161448882836);
\draw[color=ududff] (-0.32250539514428594,-1.8474432732210166) node {$A$};
\draw[color=ududff] (2.0837577938720586,0.01581362092880638) node {$B$};
\draw[color=qqqqff] (0.2205008997222343,2.4220768099451493) node {$C$};
\draw[color=qqqqff] (-2.18576228929411,0.5588199157953262) node {$D$};
\end{tikzpicture}
\end{Ill}
}

\begin{Rqs}
\begin{itemize}
\item Un carré est un \textbf{rectangle particulier}, c'est un rectangle qui a ses 4 côtés de même longueur.
\item Un carré est un \textbf{losange particulier}, c'est un losange qui a  4 angles droits.
\end{itemize}
\end{Rqs}

\section{Cercles}

\mini{.58\linewidth}{\begin{Def}
Le cercle de centre $O$ et de rayon $r$ est l'ensemble des points situés à une distance de $O$ égale à $r$.
\end{Def}
\begin{Nt}
\[M\in C(O;r)\Leftrightarrow OM=r\]
\end{Nt}
}
{.38\linewidth}{
\begin{Illqr}{Animation cercle}
\centering
\begin{tikzpicture}[line cap=round,line join=round,>=triangle 45,x=0.8cm,y=0.8cm]
\clip(-3,-3) rectangle (3,3);
\draw [line width=2.pt] (0.,0.) circle (2);
\draw [line width=1.pt,dash pattern=on 1pt off 3pt] (0.,0.)-- (2.,0.);
\draw [color=qqqqff] (0.,0.)-- ++(-2.5pt,-2.5pt) -- ++(5.0pt,5.0pt) ++(-5.0pt,0) -- ++(5.0pt,-5.0pt);
\draw[color=xdxdff] (0.07144034819024836,.3) node {$O$};
\draw [color=qqqqff] (2.,0.)-- ++(-2.5pt,0 pt) -- ++(5.0pt,0 pt);% ++(-2.5pt,-2.5pt) -- ++(0 pt,5.0pt);
\draw[color=xdxdff] (2.2,0.19681571921764635) node {$M$};
\draw[color=black] (-0.9719835125336532,1.5809494120146579) node {$C$};
\draw[color=black] (1.0296867508958722,-0.06936383708947125) node {$r$};
\end{tikzpicture}
\end{Illqr}
\begin{flushright}
Scan le QR code pour voir l'animation.
\end{flushright}
}

\begin{Voc}
\mini{.58\linewidth}{
\begin{itemize}
    \item Les points $A$, $B$, $C$, $D$ et $M$ appartiennent au cercle $\mathcal{C}$ de centre $O$.
    \item Le segment $[OM]$ est aussi appelé le \textcolor{qqqqff}{\textbf{rayon}} du cercle $\mathcal{C}$.
    \item Le segment $[CD]$ est une \textcolor{xdxdff}{\textbf{corde}} du cercle $\mathcal{C}$.
    \item La corde $[AB]$ est un \textcolor{qqwuqq}{\textbf{diamètre}} du cercle $\mathcal{C}$ car le centre O est le \textbf{milieu} de la corde.
    \item Le diamètre $[AB]$ est le double du rayon $[OM]$ : $AB=2\times OM$
    \item La portion colorée en rouge du cercle $\mathcal{C}$ est l'\textcolor{ffqqqq}{\textbf{arc de cercle}} \t{CD}
\end{itemize}
}{.38\linewidth}{
\begin{Ill}
\centering
\begin{tikzpicture}[line cap=round,line join=round,>=triangle 45,x=0.8cm,y=0.8cm]
\clip(-3,-3) rectangle (3,3);
\draw [line width=2.pt] (0.,0.) circle (2);
\draw [line width=2.pt,color=qqqqff] (0.,0.)-- (2.,0.);
\draw [line width=2.pt,color=qqwuqq] (-1.5844956505091685,1.2203989239250939)-- (0.,0.);
\draw [line width=2.pt,color=qqwuqq] (-0.7743544134251696,0.677052966073721) -- (-0.8523172720063054,0.5758304822601795);
\draw [line width=2.pt,color=qqwuqq] (-0.7321783785028615,0.6445684416649147) -- (-0.8101412370839972,0.5433459578513732);
\draw [line width=2.pt,color=qqwuqq] (0.,0.)-- (1.5844956505091685,-1.2203989239250939);
\draw [line width=2.pt,color=qqwuqq] (0.8101412370839972,-0.5433459578513732) -- (0.7321783785028615,-0.6445684416649147);
\draw [line width=2.pt,color=qqwuqq] (0.8523172720063054,-0.5758304822601795) -- (0.7743544134251696,-0.677052966073721);
\draw [line width=2.pt,color=xdxdff] (-1.9584617055069056,0.4054969149857767)-- (-0.7472818915155957,-1.855146833706937);
\draw [shift={(0.,0.)},line width=2.pt,color=ffqqqq]  plot[domain=2.9374287928705773:4.32945781065304,variable=\t]({1.*2.*cos(\t r)+0.*2.*sin(\t r)},{0.*2.*cos(\t r)+1.*2.*sin(\t r)});
\draw [color=qqqqff] (0.,0.)-- ++(-2.5pt,-2.5pt) -- ++(5.0pt,5.0pt) ++(-5.0pt,0) -- ++(5.0pt,-5.0pt);
\draw[color=qqqqff] (0.07144034819024836,0.19681571921764635) node {$O$};
\draw [color=qqqqff] (2.,0.)-- ++(-2.5pt,0 pt) -- ++(5.0pt,0 pt) ++(-2.5pt,-2.5pt) -- ++(0 pt,5.0pt);
\draw[color=qqqqff] (2.073110611619774,0.19681571921764635) node {$M$};
\draw[color=qqqqff] (1.0296867508958722,-0.06936383708947125) node {$r$};
\draw [color=qqqqff] (-1.5844956505091685,1.2203989239250939)-- ++(-2.5pt,-2.5pt) -- ++(5.0pt,5.0pt) ++(-5.0pt,0) -- ++(5.0pt,-5.0pt);
\draw[color=qqqqff] (-1.5149898074001733,1.4212416782303874) node {$A$};
\draw [color=qqqqff] (1.5844956505091685,-1.2203989239250939)-- ++(-2.5pt,-2.5pt) -- ++(5.0pt,5.0pt) ++(-5.0pt,0) -- ++(5.0pt,-5.0pt);
\draw[color=qqqqff] (1.6578705037806702,-1.0276102397950946) node {$B$};
\draw [color=qqqqff] (-1.9584617055069056,0.4054969149857767)-- ++(-2.5pt,0 pt) -- ++(5.0pt,0 pt) ++(-2.5pt,-2.5pt) -- ++(0 pt,5.0pt);
\draw[color=qqqqff] (-1.8876411862301383,0.6014086448044651) node {$C$};
\draw [color=qqqqff] (-0.7472818915155957,-1.855146833706937)-- ++(-2.5pt,0 pt) -- ++(5.0pt,0 pt) ++(-2.5pt,-2.5pt) -- ++(0 pt,5.0pt);
\draw[color=qqqqff] (-0.6738624094696813,-1.655793992679892) node {$D$};
\end{tikzpicture}
\end{Ill}
}
\end{Voc}

%\section{Les savoir-faire du parcours}
%
%\begin{CpsCol}
%\begin{itemize}
%\item Savoir décrire un polygone.
%\item Savoir reconnaître un polygone.
%\item Savoir coder une figure.                  
%\item Savoir reconnaître un polygone.
%\item Savoir construire un polygone.
%\item Savoir construire un cercle.
%\item Savoir comparer des distances avec un cercle.
%\end{itemize}
%\end{CpsCol}
%
%
%\end{document}

\end{pageCours}
